\scanpage{6}

% The torn edge cuts across the ends of the first three surviving lines.
$v_A:K^*\longrightarrow\Gamma_A$ est s\unsure{urje}\illegible{1}

Soit maintenant $v=K^*\longrightarrow\Gamma$ une valuati\unsure{on}\illegible{4} l'ensemble des $x\in K$ tels que $v(x)\geqslant0$. Si $x\in A$ et $y\in$ \illegible{4} $xy\in A$, $-x\in A$ ; en outre, $1\in A$ ; donc $A$ est un sous-anneau de $K$ (contenant 1). Si $x\in K$ et $x\notin A$, on a $v(x)<0$, d'où $v(x^{-1})>0$ et $x^{-1}\in A$ Donc $A$ est un anneau de valuation pour $K$. Les élément inversibles de $A$ forment un groupe $U$ qui est l'ensemble des $x\in K^*$ pour lesquels $v(x)=0$. On peut donc factoriser $v$ sous la forme $v=\sigma\mathbin{\unsure{\mathrm{o}|\circ}}v_A$, où $\sigma$ est un homomorphisme injectif de $\Gamma_A$ dans $\Gamma$. Si $\xi=(x)$ est un élément $\geqslant0$ de $\Gamma_A$, on a $x\in A$, d'où $\sigma(x)=v(x)\geqslant0$ ; donc $\sigma$ est croissant, ce qui achève de prouver la prop.3 .

La factorisation $v=\sigma\circ v_A$ est appelée la \uline{factorisation canonique} de $v$. On dit que $A$ est \uline{l'anneau de la valuation} $v$ (d'où la terminologie d'anneau de valuation), que $\mathfrak{m}(A)$ est \uline{l'idéal de} $v$, et que $\unsure{\varkappa|\mathfrak{k}|x}_A$ est le \uline{corps résiduel de} $v$. Le sous-groupe $\sigma(\Gamma_A)=v(K^*)$ est appelé le \uline{groupe des ordres de} $v$. Pour $x\in K$, l'élément $v(x)$ est parfois appelé la valuation ou l'ordre de $x$ (A vider ! ! ). Deux valuations $v$ et $v'$ sont dites \uline{équivalentes} si elles ont même anneau de valuation ; pour cela, il faut et il suffit, en vertu de la prop.3, qu'il existe un isomorphisme $\lambda$ du groupe ordonné $v(K^*)$ sur le groupe ordonné $v'(K^*)$ tel que $v'=\lambda\circ v$.

\uline{Proposition} 4 - \uline{Soit} $v$ \uline{une valuation d'un corps} $K$. \uline{Quels que soient} $x_i\in K$ $(1\leqslant i\leqslant n)$, \uline{on a}
\[
v(\sum\nolimits_{i=1}^n x_i)\geqslant\inf\nolimits_{1\leqslant i\leqslant n}(v(x_i))\tag*{(5)}
\]

Si en outre il existe un seul indice $k$ tel que $v(x_k)=\inf_{1\leqslant i\leqslant n}(v(x_{\unsure{i}}$
% right edge cuts off the end of this line
il y a alors égalité dans (5). En particulier, si $v(x)\ne v(y)$, on a

% archive sheet number: 5

